\documentclass[11pt]{article}
\usepackage{geometry}
\geometry{letterpaper, margin=1in}
\usepackage{graphicx}
\usepackage{listings}
\usepackage{xcolor}
\usepackage[utf8]{inputenc}
\usepackage[T1]{fontenc}
\usepackage{lmodern}
\usepackage{float}
\usepackage{booktabs}
\usepackage{tabularx}
\usepackage{hyperref}

\lstset{
    basicstyle=\ttfamily\small,
    breaklines=true,
    frame=single,
    showstringspaces=false
}

% Screenshots live in ./screenshots
\graphicspath{{screenshots/}}

% \screenshot{file}{caption}{label}
% Shows screenshots/<file> if it exists; otherwise a placeholder box so the
% report still compiles before the screenshots are added.
\newcommand{\screenshot}[3]{%
  \begin{figure}[H]
    \centering
    \IfFileExists{screenshots/#1}{%
      \includegraphics[width=\textwidth,height=0.8\textheight,keepaspectratio]{#1}%
    }{%
      \fbox{\parbox[c][4cm][c]{0.9\textwidth}{\centering
        \textbf{Screenshot missing}\\[4pt]
        Save it as \texttt{report/screenshots/\detokenize{#1}}}}%
    }
    \caption{#2}
    \label{#3}
  \end{figure}}

\newcommand{\question}[2]{\subsection*{Question #1}\noindent\textit{#2}\par\medskip}
\newcommand{\code}[1]{\texttt{#1}}

\begin{document}
\sloppy

\title{CS-375 Operating Systems Lab 9:\\Advanced Segmented, Paged Memory Model Simulator\\
with TLB, Multi-Level Paging, and Dynamic Memory Management}
\author{Santiago Soto}
\date{\today}
\maketitle

\tableofcontents
\newpage

%=====================================================================
\section{Introduction}
This report documents a simulator of a segmented, paged virtual memory system.
It models a segment table, a two-level page table, a TLB with LRU replacement,
demand paging with LRU or FIFO page replacement, and segments created and
removed at runtime. The simulator is in
\code{memory\_simulator\_advanced.cpp}, and the sample inputs are in
\code{segments.txt} and \code{addresses.txt}.

To build and run the simulator:
\begin{lstlisting}
g++ -o memory_simulator_advanced memory_simulator_advanced.cpp -std=c++11
./memory_simulator_advanced
\end{lstlisting}

Every result below uses the option \code{--seed}, so the same command repeats
the same run. Experiments use 1000 addresses unless stated otherwise.

%=====================================================================
\section{Design}

\subsection{Data Structures}
\begin{table}[H]
\centering
\begin{tabularx}{\textwidth}{lX}
\toprule
\textbf{Structure} & \textbf{Role} \\
\midrule
\code{Segment} & ID, base address, limit in pages, protection and per-segment counters. \\
\code{Page} & Frame number, present bit, protection and \code{last\_access}. \\
\code{PageTable} & Second level: \code{getFrameNumber}, \code{setFrame} and \code{replacePage}. \\
\code{DirectoryTable} & First level: creates each page table the first time it is used. \\
\code{TLB} & Fixed-size cache of \code{seg:dir:page} to frame, with LRU replacement. \\
\code{PhysicalMemory} & Frame allocation, the FIFO queue, the LRU list and the owner of each frame. \\
\code{SegmentTable} & Coordinates the above and implements \code{translateAddress}, metrics, the timeline and the log. \\
\bottomrule
\end{tabularx}
\caption{Main data structures}
\end{table}

A segment's limit is its total number of pages. Each second-level table holds 4
pages, so a 16-page segment has directories 0--3 and a 5-page segment has
directories 0--1, the second holding a single page.

\subsection{Address Translation}
\code{translateAddress(segNum, dirNum, pageNum, offset, accessType)} works as follows:
\begin{enumerate}
  \item Draw the base latency (random, 1--10 units).
  \item Validate the segment, segment protection, directory, page and offset.
        A failure raises an exception and is recorded with its reason.
  \item Check the page's protection. This runs on every access, including TLB hits.
  \item Look up the TLB. A hit updates \code{last\_access} and finishes.
  \item On a miss, walk both table levels, adding 10 units per level.
  \item If the page is not present, \code{getFrameNumber} raises
        \code{PageFaultException}. \code{translateAddress} catches it, loads the
        page with \code{replacePage}, invalidates the victim's TLB entry and adds
        100 units. This mirrors a hardware trap followed by a restart.
  \item Update the TLB and metrics, then return
        $\mathit{base} + \mathit{frame} \times \mathit{pageSize} + \mathit{offset}$.
\end{enumerate}
Bounds are checked before the TLB lookup, as in the starter code, so invalid
addresses do not count as TLB misses.

\subsection{Corrections to the Starter Code}
The starter code did not compile, and several of its defects would have
invalidated the experiments. Each correction has a regression test in
\code{tests/test\_simulator.cpp}.

\begin{table}[H]
\centering
\small
\begin{tabularx}{\textwidth}{XX}
\toprule
\textbf{Starter code problem} & \textbf{Correction} \\
\midrule
Missing \code{<queue>}, \code{<algorithm>} and \code{<chrono>}; \code{segments} was private but used outside the class. & Added the includes and accessor methods. \\
Demand page faults were never counted, so LRU and FIFO could not be compared. & \code{fault\_count} now counts demand faults; invalid addresses and violations have their own counters. \\
An evicted page stayed present, so two pages could share one frame, and its TLB entry stayed valid. & Each frame records its owner; eviction clears that page and its TLB entry. \\
Page tables were filled with random frames that were never allocated. & Pages only hold frames obtained from \code{PhysicalMemory}. \\
TLB hits skipped the page protection check and the \code{last\_access} update. & Both now happen on every access. \\
A page fault overwrote the page's protection, so a write fault made a read-only page writable. & A page fault no longer changes protection. \\
\code{--pagesize} was ignored (1000 was hard-coded). & The configured page size is used throughout. \\
Segments were indexed by position, so \code{add 3 ...} could not be reached or removed. & Segments are stored in a map keyed by ID. \\
An invalid segment caused an out-of-bounds write; a negative \code{dirNum} never finished allocating. & All values are validated, with detailed messages such as ``Invalid dirNum 5, max 3''. \\
\code{removeSegment} left stale TLB entries and leaked its tables. & It invalidates the TLB and frees every table; \code{DirectoryTable} has a destructor. \\
\code{TLB::put} could evict an entry or add a duplicate when updating a key. & Updates replace the entry in place. \\
Invalid addresses were produced by multiplying by 0, which yields valid segment, directory and page 0. & Each invalid address now breaks exactly one rule. \\
LRU used list order, not the required timestamps; time only advanced on TLB misses. & LRU evicts the smallest \code{last\_access}; time advances on every translation. \\
\bottomrule
\end{tabularx}
\caption{Defects corrected in the starter code}
\end{table}

Three design decisions keep or deliberately change the starter's behavior:
\begin{itemize}
  \item A segment's limit is its total number of pages, as the lab defines it. In
        the starter code each directory had \code{limit} pages and the number of
        directories had no bound, so \code{dirNum} could never be invalid.
  \item The physical address formula is unchanged, including the segment base.
        As a result, two different frames in different segments can produce the
        same address number. This does not affect any metric.
  \item The random 1--10 latency is kept, with fixed costs added for table walks
        and page faults. Without them the TLB could have no effect on latency.
        These costs are a model; a real disk access costs far more than 100
        times a memory access.
\end{itemize}

%=====================================================================
\section{Test Output}

\subsection{Compilation and Regression Tests}
The program compiles with no warnings under \code{-Wall -Wextra}. All 73 checks
pass with AddressSanitizer and UndefinedBehaviorSanitizer enabled, and Valgrind
reports no leaks or memory errors.

\screenshot{01_compile.png}{Compilation and regression tests}{fig:compile}

\subsection{Manual Translations}
All manual runs use \code{segments.txt}, in which pages 0, 1 and 5 of segment 0
and page 0 of segment 1 start in frames 0--3.

\begin{table}[H]
\centering
\small
\begin{tabular}{lll}
\toprule
\textbf{Address (seg dir page offset access)} & \textbf{Result} & \textbf{Reason} \\
\midrule
\code{0 0 1 250 0} & 1250 & Resident page, TLB miss \\
\code{0 0 1 300 1} & 1300 & TLB hit, write allowed \\
\code{0 1 3 10 0} & 4010 & Page fault, loaded into frame 4 \\
\code{1 0 0 0 1} & Protection Violation & Write to read-only segment \\
\code{5 0 0 0 0} & Segmentation Fault & Segment 5 does not exist \\
\code{0 7 0 0 0} & Invalid dirNum & Segment 0 has directories 0--3 \\
\code{0 0 0 1500 0} & Offset Fault & Offset exceeds page size 1000 \\
\bottomrule
\end{tabular}
\caption{Manual translations}
\end{table}

\screenshot{02_manual_translations.png}{Manual translations, valid and invalid}{fig:manual}

\subsection{Batch Processing}
\code{addresses.txt} contains 21 addresses, run with 6 frames and a 4-entry TLB.
Of these, 15 succeeded and 6 failed: 1 segmentation fault, 3 invalid addresses
and 2 protection violations. There were 8 page faults and 6 evictions.

\screenshot{03_batch_results.png}{Contents of batch\_results.txt}{fig:batch}

\subsection{Random Address Test}
With 1000 addresses, a valid ratio of 0.7, 6 frames and a 4-entry TLB, 691
translations succeeded and 309 failed. There were 170 page faults and a TLB hit
rate of 68.74\%. The log gives a detailed reason for each failure.

\screenshot{04_random_results.png}{Start and end of random\_results.txt}{fig:random}

\subsection{Dynamic Segments}
Adding segment 3 and touching two of its pages raised utilization from 4 to 6
of 12 frames, and each of its two second-level tables was allocated on first
use. Removing the segment returned utilization to 4 of 12 frames (33.33\%) and
removed its TLB entries. A later access to segment 3 raised a segmentation
fault.

\screenshot{05_add_map.png}{Memory map after adding segment 3}{fig:addmap}
\screenshot{06_add_stats.png}{Statistics after adding segment 3}{fig:addstats}
\screenshot{07_remove_map.png}{Memory map after removing segment 3}{fig:removemap}
\screenshot{08_remove_stats.png}{Statistics after removing segment 3}{fig:removestats}

%=====================================================================
\section{Experiments}
Each comparison replays one seeded address sequence on a new simulator with
memory initially empty, using \code{segments.txt} (40 pages in total). The
\emph{locality} of a workload is the fraction of references sent to a hot region
of 20\% of one segment, which moves every 100 references. Tables show seed 42;
\code{scripts/seed\_sweep.sh} repeats each experiment over seeds 1--10.

\subsection{LRU vs FIFO}
\begin{table}[H]
\centering
\begin{tabular}{lrrrrr}
\toprule
\textbf{Workload} & \textbf{Frames} & \textbf{Policy} & \textbf{Page faults} & \textbf{TLB hit \%} & \textbf{Avg. latency} \\
\midrule
Locality 0.8 & 12 & LRU & 158 & 81.50 & 25.12 \\
             & 12 & FIFO & 181 & 78.90 & 27.94 \\
Locality 0.8 & 6 & LRU & 212 & 78.80 & 31.06 \\
             & 6 & FIFO & 282 & 71.80 & 39.46 \\
Uniform (0.0) & 12 & LRU & 696 & 20.50 & 91.12 \\
             & 12 & FIFO & 682 & 20.30 & 89.76 \\
\bottomrule
\end{tabular}
\caption{LRU vs FIFO, 8-entry TLB, all addresses valid}
\end{table}

Over 10 seeds with 12 frames, LRU averaged 168.7 faults against 201.0 for FIFO
at locality 0.8, and had fewer faults on all 10 seeds. On the uniform workload
the averages were 678.9 and 678.3, and each policy had fewer faults on 5 seeds.

\screenshot{09_lru_vs_fifo.png}{LRU vs FIFO on the same workload}{fig:policy}

\subsection{TLB Size}
\begin{table}[H]
\centering
\begin{tabular}{lrrrr}
\toprule
 & \multicolumn{2}{c}{\textbf{Seed 42}} & \multicolumn{2}{c}{\textbf{Mean of 10 seeds}} \\
\textbf{TLB} & \textbf{Hit \%} & \textbf{Avg. latency} & \textbf{Hit \%} & \textbf{Avg. latency} \\
\midrule
Off & 0.00 & 36.72 & 0.00 & 37.62 \\
4 entries & 70.50 & 22.62 & 66.94 & 24.24 \\
16 entries & 86.70 & 19.38 & 85.59 & 20.51 \\
\bottomrule
\end{tabular}
\caption{TLB size, 20 frames, LRU, locality 0.8, 111 page faults at seed 42 for every size}
\end{table}

\screenshot{10_tlb_size.png}{TLB size comparison}{fig:tlb}

\subsection{Page Size}
Segment sizes and physical memory stay fixed in bytes, at 40\,000 bytes of
segments and 20\,000 bytes of memory, so only the page size changes.

\begin{table}[H]
\centering
\begin{tabular}{lrrrrr}
\toprule
\textbf{Page size} & \textbf{Frames} & \textbf{Page faults} & \textbf{TLB hit \%} & \textbf{Avg. latency} & \textbf{Page table entries} \\
\midrule
1000 & 20 & 111 & 81.50 & 20.42 & 50 \\
4000 & 5 & 101 & 89.90 & 17.74 & 13 \\
\bottomrule
\end{tabular}
\caption{Page size with equal memory, LRU, 8-entry TLB, locality 0.8}
\end{table}

\screenshot{11_page_size.png}{Page size comparison}{fig:pagesize}

\subsection{Valid/Invalid Ratio}
\begin{table}[H]
\centering
\begin{tabular}{rrrrrr}
\toprule
\textbf{Valid ratio} & \textbf{Failed} & \textbf{Page faults} & \textbf{Fault \% of successes} & \textbf{TLB hit \%} & \textbf{Avg. latency} \\
\midrule
1.0 & 0 & 158 & 15.80 & 81.50 & 25.12 \\
0.9 & 99 & 144 & 15.98 & 80.80 & 23.48 \\
0.7 & 309 & 125 & 18.09 & 78.58 & 21.08 \\
0.5 & 506 & 99 & 20.04 & 77.73 & 17.72 \\
0.1 & 893 & 35 & 32.71 & 63.55 & 9.90 \\
\bottomrule
\end{tabular}
\caption{Valid ratio, 12 frames, LRU, 8-entry TLB, locality 0.8}
\end{table}

\screenshot{12_valid_ratio.png}{Valid/invalid ratio comparison}{fig:ratio}

%=====================================================================
\section{Short-Answer Questions}

\question{1}{What is the role of the TLB in reducing address translation overhead?}
The TLB caches recent page-to-frame translations, so a hit avoids walking the
page tables. In the simulator a hit saves 20 latency units, and at a 68.74\% hit
rate the hits avoided 26.07\% of the total translation cost.

\question{2}{How does a two-level page table differ from a single-level page table in structure?}
A single-level table is one array with an entry for every page, while a
two-level table uses a directory that points to smaller page tables. The
second-level tables are created only when needed, so unused address ranges use
no table memory.

\question{3}{Why is LRU often preferred over FIFO for page replacement in modern systems?}
LRU uses recent use as a predictor of future use, so it keeps a program's
working set in memory, while FIFO can evict a heavily used page just because it
was loaded long ago. LRU had fewer page faults on all 10 seeds of the workload
with locality.

\question{4}{What triggers a segmentation fault vs. a page fault in your simulator?}
A segmentation fault is raised when the segment number does not match an existing
segment, and it is fatal to that translation. A page fault is raised when a valid
page is not present, and it is resolved by loading the page and completing the
translation.

\question{5}{How does the \code{present} bit in the page table affect translation?}
If the bit is set, the stored frame number is used directly. If it is clear,
\code{getFrameNumber} raises a page fault so the page can be loaded, possibly
evicting another page.

\question{6}{Why is memory cleanup critical when removing a segment?}
Frames that are not freed are lost to every other segment, and stale TLB entries
could translate a later segment with the same ID to frames it does not own.
Table memory must also be deleted, or the simulator itself leaks memory.

\question{7}{What is the purpose of tracking \code{last\_access} in the page table?}
It records when each page was last used, so LRU can evict the resident page with
the smallest timestamp. It must be updated on TLB hits too, or frequently used
pages would look old.

\question{8}{How does the protection flag in segments enhance system security?}
It lets the system enforce read-only access to regions such as code or shared
data, so a buggy or malicious write is stopped before it can change memory. The
simulator checks it on every access and logs each violation.

\question{9}{Why might a smaller page size increase page faults?}
The same data spans more pages, so a program must fault in more pages to touch
it. With equal memory, a page size of 1000 caused 111 faults against 101 with a
page size of 4000.

\question{10}{What is the benefit of logging page fault statistics per segment?}
It shows which segment is causing memory pressure, rather than a single overall
rate that hides the cause. That lets memory be directed to the segments that need
it, and the simulator's suggestions are based on these per-segment rates.

%=====================================================================
\section{Long-Answer Questions}

\question{11}{Explain how the TLB improves performance in your simulator. Use your code's metrics (e.g., hit rate) to discuss its effectiveness in a scenario with frequent address translations.}
% BEGIN Q11
The TLB is a small cache in front of the two-level page table. Without it, every
translation must read a directory entry and then a page table entry, which the
simulator charges at 10 latency units per level. On a hit, \code{translateAddress}
takes the frame directly from the cache, skips the walk and only pays the base
latency of 1--10 units. Protection and \code{last\_access} are still handled, so
the shortcut does not change behavior, only cost.

The random address test shows the size of the effect. With 1000 addresses, 6
frames and a 4-entry TLB, the hit rate was 68.74\% (475 hits and 216 misses). The
average latency was 5.75 units for a hit, 24.70 for a miss that needed only a
table walk, and 125.75 for a page fault. The simulator reports that the hits
avoided 9500 units of table walks, or 26.07\% of what the run would have cost
without a TLB.

The comparison table isolates the TLB by replaying one workload with its size as
the only difference. With 20 frames, every run had 111 page faults, because a TLB
does not change which pages are resident. Average latency still fell from 36.72
units with the TLB off to 22.62 with 4 entries and 19.38 with 16 entries. That is
a 47\% reduction with 16 entries, and it came entirely from avoided walks.

The result depends on locality. A TLB works because programs reuse a small set of
pages, so a few entries cover most references. On the uniform workload, with 12
frames and 8 entries, the hit rate fell to about 20\%. Frequent translation alone
does not make a TLB effective: the references must also repeat. Because page
faults dominate latency in this model, the TLB improves the common case, while
the page fault rate still sets the worst case.
% END Q11

\question{12}{Compare LRU and FIFO page replacement in your simulator. Run tests with both policies on the same workload and analyze differences in page fault rates and performance.}
% BEGIN Q12
Both policies were tested with \code{--compare policy}, which replays one seeded
address sequence through a new simulator for each policy. Every address was
valid, so all differences come from replacement decisions.

On a workload with locality 0.8, 12 frames and an 8-entry TLB, LRU had 158 page
faults (15.8\%) and FIFO had 181 (18.1\%). Average latency was 25.12 units for LRU
and 27.94 for FIFO, because each fault adds 100 units. The TLB hit rate also
improved, from 78.9\% to 81.5\%, since eviction removes the victim's TLB entry and
fewer evictions leave more useful entries in place. Reducing memory to 6 frames
widened the gap: LRU had 212 faults and FIFO had 282, which is 33\% more. The less
memory there is, the more each wrong eviction costs.

One seed could be a coincidence, so the test was repeated with seeds 1--10. LRU
averaged 168.7 faults against 201.0 for FIFO and had fewer faults on all 10 seeds.
The reason is that FIFO looks only at load time. A hot page is evicted once it
becomes the oldest, however often it is used, and then it immediately faults back
in. LRU keeps such a page because every use refreshes its \code{last\_access}.

The advantage disappears without locality. On a uniform workload LRU had 696
faults and FIFO had 682, and over 10 seeds the averages were 678.9 and 678.3, with
each policy ahead on 5 seeds. When the past does not predict the future, recency
provides no information and the two policies are equivalent. Near 70\% of accesses
fault here because 12 frames hold only 30\% of the 40 pages.

LRU also costs more. FIFO only updates a queue when a page loads, while LRU
updates a timestamp on every access and scans frames to find the minimum. Real
systems therefore approximate LRU with schemes such as the clock algorithm.
% END Q12

\question{13}{Discuss the challenges of implementing two-level paging. How does it impact memory usage and translation latency compared to single-level paging?}
% BEGIN Q13
The main design challenge was deciding what a segment's limit means when an
address has separate directory and page numbers. In the starter code each
directory had \code{limit} pages and the number of directories was unbounded, so
any \code{dirNum} allocated tables, and a negative value never finished. In this
simulator the limit is the total page count, and each table holds 4 pages. That
makes both fields checkable, and it allows messages such as ``Invalid dirNum 5,
max 3''. The last table may be partial: a 5-page segment has one page in
directory 1, so bounds must be checked against the actual table.

Other challenges concern ownership. Tables are created lazily, so a directory
entry may be null and every traversal must handle that. Cleanup must free each
table exactly once. Eviction also needs to reach back to a page from its frame,
which is why each frame records its owning segment, directory and page.

Memory usage depends on how sparsely a segment is used. A single-level table needs
one entry per page, so the three segments in \code{segments.txt} would always
need 40 entries. A two-level design starts with only 10 directory entries. After
the dynamic test it held 22 entries, against 40 for a single level. It is not
always smaller: in the 1000-address stress test every table was eventually
allocated, giving 50 entries against 40, because the directories are extra. Two
levels help sparse address spaces, while a densely used one pays a small overhead.
With real 32-bit or 64-bit address spaces, which are overwhelmingly sparse, the
saving is very large.

The cost is latency. A TLB miss needs two memory references instead of one, so the
walk costs 20 units here rather than 10. With the TLB, this matters only on
misses. At the 18.5\% miss rate measured with an 8-entry TLB, the extra level adds
about 1.85 units per access, against an average of 25.12. Without a TLB it would
add a full 10 units to every access. This is why multi-level paging and TLBs are
used together.
% END Q13

\question{14}{Describe how your simulator handles dynamic segment creation and deletion. What measures prevent memory leaks or invalid accesses?}
% BEGIN Q14
Segments are created with \code{add <id> <base> <limit> <RO|RW>}. \code{addSegment}
rejects a duplicate ID, a negative ID or base, and a limit outside 1--65536, then
stores a \code{Segment} in a map keyed by ID and creates a \code{DirectoryTable}.
The map matters: the starter code indexed a vector by position, so \code{add 3}
on a two-segment table stored the segment at index 2 and could never remove it.
No frames or second-level tables are allocated at creation. They appear on first
use, so adding a segment costs only its directory.

Deletion is handled by \code{removeSegment}, in a specific order. It first walks
every allocated table and returns each present page's frame through
\code{freeFrame}, which also removes the frame from the FIFO queue or LRU list and
clears its owner. If a frame stayed in those structures, a later replacement could
choose a frame belonging to nobody. Next it calls \code{invalidateSegment}, which
removes all of that segment's TLB entries. Finally it deletes the directory and
erases the segment.

Several measures prevent leaks. \code{DirectoryTable} has a destructor that frees
its page tables, a gap in the starter code, and its copy operations are disabled
so that tables cannot be freed twice. Valgrind reports that all heap blocks were
freed after a session that adds a segment, runs a stress test and removes it.

Several measures prevent invalid accesses. Each translation looks the segment up
in the map, so a removed segment raises a segmentation fault instead of reading a
blank entry. In the test, \code{translate 3 0 0 100 0} after removal did exactly
this. TLB invalidation stops a stale hit, which would otherwise be the most
dangerous failure: a new segment reusing the ID could reach another segment's
frames with no error at all. A regression test removes a segment, recreates it
with the same ID, and confirms that its first access faults. The memory map also
confirms cleanup: after removal, utilization returned from 6 to 4 of 12 frames.
% END Q14

\question{15}{Analyze the impact of TLB size on performance. Test with sizes 4 and 16, and discuss trade-offs between hit rate and memory overhead.}
% BEGIN Q15
TLB sizes were tested with \code{--compare tlb} on one workload of 1000 valid
addresses with locality 0.8, 20 frames and LRU replacement. A disabled TLB was
included as a baseline. Over seeds 1--10, the mean hit rate rose from 66.94\% with
4 entries to 85.59\% with 16, and mean latency fell from 37.62 units with no TLB to
24.24 and then 20.51. Seed 42 behaved the same way: 70.50\% and 86.70\% hit rates,
with latencies of 36.72, 22.62 and 19.38. The page fault count was 111 at seed 42
for all sizes, since the TLB affects only how quickly resident pages are found.

The results show diminishing returns. The first 4 entries saved 13.38 units per
access on average. The next 12 entries saved 3.73 units, so each added entry is
worth about a tenth of what each of the first four was worth. This follows from
the workload: the hot region covers only about 20\% of a segment, or roughly 3 to 4
pages, so four entries capture much of it. The remaining entries mostly help with
the 20\% of references outside that region and with moves of the hot region.

The size that helps also has a natural ceiling. A TLB entry is useful only for a
resident page, so with 20 frames a TLB of more than 20 entries could not raise the
hit rate here, and a smaller physical memory lowers that limit.

The memory overhead in the simulator is trivial, but in hardware it is the real
constraint. A TLB is fully or highly associative, so a lookup compares many entries
at once. More entries take more chip area and power and can lengthen the lookup,
which runs on every memory reference. A simulator with a fixed hit cost cannot show
that penalty. On these numbers, 16 entries is justified because it adds nearly 19
percentage points of hit rate. If each extra entry were expensive, though, 4
entries would already deliver most of the benefit. Real designs resolve this with
a small, fast first-level TLB backed by a larger, slower second level.
% END Q15

\question{16}{Explain how protection enforcement in your simulator mimics real OS behavior. Provide an example of a protection violation and its handling.}
% BEGIN Q16
Real systems enforce protection in hardware on every access. The memory management
unit compares the access type with permission bits in the page table entry, which
the TLB also caches. A denied access traps to the kernel, and the instruction does
not complete. On Linux the process normally receives \code{SIGSEGV}. Protection is
therefore not an optional check that fast paths may skip.

The simulator follows this model at two levels. Each segment has a protection flag
and each page has its own, so a page may be stricter than its segment but can
never be more permissive. For example, page 2 of segment 0 in \code{segments.txt}
is read-only inside a read-write segment. A \code{ProtectionViolation} exception
plays the role of the trap: it stops the translation, so no address is returned,
no page is loaded and no frame changes. The check runs before the page fault
handler, so a forbidden write cannot bring a page into memory.

Two starter defects would have broken this. A TLB hit returned early, so a write
could pass after a read cached the translation. A real TLB holds permission bits
for this reason, and here the page's protection is rechecked on each access. A
page fault also copied the access type into the page's protection, so a write
fault turned a read-only page into a writable one.

As an example, \code{translate 1 0 0 0 1} is a write to segment 1, which is
read-only. The simulator performs these steps:
\begin{enumerate}
  \item It finds the segment and sees a write to a read-only region.
  \item It raises ``Protection Violation: Cannot write to read-only segment 1''.
  \item The handler in \code{translateAddress} increments the segment's violation
        counter and the global counter.
  \item The event is added to the timeline and \code{memory\_log.txt}, with its
        time, full logical address and latency.
  \item The exception reaches the command loop, which prints the error and continues.
\end{enumerate}

The last step is where the simulator differs from a real OS, which would normally
terminate the process. The simulator continues so a batch can report every result,
and each violation is logged with its context, much as a kernel records a
fault.
% END Q16

\question{17}{Evaluate the effectiveness of your stress testing with random addresses. How does varying the valid/invalid ratio reveal system strengths and weaknesses?}
% BEGIN Q17
Each stress address is either fully valid or breaks exactly one rule, chosen from
five: segment, directory, page, offset, or a write to a read-only segment. This
replaced the starter's generator, which multiplied values by a Boolean. That
produced 0, a valid value, when an address was meant to be invalid, and it compared
\code{mt19937} output with \code{RAND\_MAX}, which has a different range. The test
requires a ratio of 1.0 to give no failures and 0.0 to give all failures, and both
pass. At seed 42, the failures were 99, 309, 506 and 893 for ratios of 0.9, 0.7, 0.5
and 0.1, close to the expected 100, 300, 500 and 900.

Varying the ratio showed clear strengths. Every invalid address was rejected with a
specific reason, and none corrupted state. Even at a ratio of 0.1, with 893 failures,
the run finished normally. Sanitizers and a test that checks that no two pages share
a frame also pass after stress runs, which is stronger evidence than the output
looking plausible.

It also revealed how metrics behave under failure. As the ratio fell from 1.0 to 0.1,
average latency dropped from 25.12 to 9.90 units. That looks like an improvement, but
failed translations are simply cheap, because they stop before any walk or fault. A
rate computed over all attempts would therefore mislead, so the simulator reports page
faults relative to successful translations. That rate rose from 15.80\% to 32.71\%,
and the TLB hit rate fell from 81.50\% to 63.55\%. Fewer valid references mean that a
larger share of them are first-time accesses, and the hot region gets too few
references to build up reuse before it moves.

The testing also has weaknesses. The generator never writes to a read-only page inside
a read-write segment, so page-level violations occur only in manual and batch runs. Faults
cost a fixed amount, and synthetic locality only approximates real programs. Passing
also shows only that the checks I wrote held, not that the model is complete: for
example, the simulator has no concurrency, so race conditions cannot appear in it.
% END Q17

\question{18}{Discuss how page fault frequency tracking can inform memory management decisions. Propose an algorithm to dynamically adjust segment limits based on your simulator's data.}
% BEGIN Q18
Page fault frequency (PFF) is a direct and inexpensive signal of whether a segment has
enough memory. A high rate means that its working set does not fit, so pages are
evicted and quickly faulted back in. A very low rate suggests that it holds more
memory than it needs. The simulator tracks faults and accesses per segment, and it
suggests a larger limit when a segment's rate passes 20\% over at least 10 accesses.
In the stress test, segments 0 and 2 reached 21.93\% and 23.02\% and received
suggestions, while segment 1 stayed at 12.47\%.

The data also shows a limitation in the lab's rule. A segment's limit is the size of
its address range, not the memory it holds, so raising a limit does not reduce demand
faults. What reduces them is more frames. Invalid-page errors are the signal that a
limit is too small. A sound algorithm should separate these cases, so the proposal
below adds a per-segment frame quota and evicts within a segment that is at its quota.

\begin{lstlisting}
Every W = 100 translations, for each segment s:
  fault_rate   = faults in window / accesses in window
  bounds_rate  = bounds errors in window / accesses in window

  if bounds_rate > 5%:
      s.limit = min(ceil(s.limit * 1.5), MAX_PAGES)
  if fault_rate > 20%:
      if a free frame exists: s.quota += 1
      else: take 1 frame from the segment with the
            lowest fault_rate, if that rate is below 5%
  else if fault_rate < 5% and s.quota > MIN_QUOTA:
      s.quota -= 1
  if total demand exceeds frames and no donor exists:
      suspend the segment with the highest fault_rate
\end{lstlisting}

The gap between 5\% and 20\% is deliberate. With a single threshold, a segment near
that value would gain and lose frames at every window. Using a window means that old
behavior ages out, and the 10-access minimum already in the simulator avoids acting on
tiny samples, such as a first access that counts as a 100\% fault rate. Cold-start
faults should be discounted for the same reason. The final step handles thrashing:
when total demand exceeds memory, taking frames from one segment only moves the
problem, so a segment is suspended instead.
% END Q18

\question{19}{Propose a real-world application (e.g., IoT device, cloud VM) for your simulator. How would you extend it to meet specific requirements (e.g., real-time deadlines, multi-core support)?}
% BEGIN Q19
A strong application is capacity planning for a cloud host running many virtual
machines. An operator must decide how much memory and TLB capacity each VM needs
before committing hardware. Instead of testing on production systems, memory traces
from each VM could be replayed through the simulator under different configurations.
The existing \code{--batch} input and \code{--compare} tables already support this
workflow, and page fault and hit rates per segment map naturally to per-VM metrics
if each VM is treated as a segment. Three extensions would be needed.

The first is multi-core support. The simulator has one TLB, whereas each core has its
own. I would give each core a \code{TLB} object and run workloads in threads, with
locks around shared frame structures. The important change is TLB shootdown: when a
page is evicted, every core that may cache it must be told to invalidate that entry,
which interrupts the other cores. Modeling that cost would show how eviction scales
with core count. Address space identifiers would also be needed, so that a context
switch does not require flushing the whole TLB. The current keys already contain a
segment number, which is a similar idea.

The second is realistic fault costing. Writing back a modified page costs much more
than dropping a clean one, so \code{Page} needs a dirty bit, set on writes, and the
replacement policy should be able to prefer clean victims. Fault latency should also
model storage more closely, as a fixed 100 units is far too low for a disk.

The third is real-time support for tenants with latency targets. Average latency is
not enough, because one page fault can miss a deadline even when the average looks
fine. The simulator would need to report worst-case and 99th-percentile latency. It
would also need page locking, with pages excluded from eviction, similar to
\code{mlock}. A deadline could be attached to each request, with misses counted
as a metric. Finally, I would replace the LRU scan, which is linear in the number of
frames, with a clock algorithm, so that it can handle the large memory sizes that
such hosts have.
% END Q19

\question{20}{Analyze the trade-offs of page size in your simulator. Test with sizes 1000 and 4000, and discuss impacts on page faults, memory utilization, and translation overhead.}
% BEGIN Q20
A fair test required care. If only \code{--pagesize} changes, page and frame counts
stay the same, so the fault count cannot change and memory simply becomes four times
larger. \code{--compare pagesize} instead keeps bytes constant: 40\,000 bytes of
segments and 20\,000 bytes of memory. That gives 20 frames at a page size of 1000 and
5 at 4000. The workload is generated as byte positions, so both runs reference the
same bytes.

Page faults fell slightly with larger pages. At seed 42 there were 111 faults at 1000
and 101 at 4000. Over seeds 1--10 the means were 121.2 and 105.4, a 13\% reduction.
A larger page brings in neighboring data, which is prefetching: one fault at 4000
covers what could take four faults at 1000. On the uniform workload the means were
478.8 and 470.8, a gain of under 2\%. With no spatial locality, neighboring bytes are
rarely needed. There is also a risk that the test understates: with only 5 frames, the
simulator can hold just 5 distinct regions, so a workload with many scattered hot
spots could fault more with large pages.

Translation overhead improved the most. Each TLB entry now covers four times as much
memory, so the reach of 8 entries grows from 8000 to 32\,000 bytes. The mean hit rate
rose from 80.04\% to 89.46\%, and on the uniform workload from 21.61\% to 52.92\%.
Mean latency fell from 21.62 to 18.15 units. Page table size also dropped from 50
entries to 13, as fewer pages need describing.

Memory utilization is where the cost appears, and here the metric hides it. Both
runs report 100\% of frames in use, but that counts frames, not useful bytes. A page
is loaded whole even if one byte is needed, so larger pages increase internal
fragmentation, and the wasted space is greater with 4000-byte pages. The simulator
does not measure this; it would need tracking of which bytes in a page are used.

Overall, 4000 is better on every metric measured here, but those metrics do not
include fragmentation. Real systems balance these factors, commonly using 4\,KB pages
with larger huge pages for workloads that need more TLB reach.
% END Q20

%=====================================================================
\section{Discussion}

\subsection{Challenges}
\begin{itemize}
  \item \textbf{Consistency after eviction.} A page's state exists in the page table,
        the frame table and the TLB. The starter code updated only some of these, which
        let two pages share a frame without any visible error. Recording each frame's
        owner solved this, and a test now checks for shared frames after stress runs.
  \item \textbf{Making experiments comparable.} The original workload used different
        random numbers in each run, and the default segments fit in memory, so no
        replacement ever happened. Seeding, a fixed draw order and a 40-page
        configuration made comparisons fair.
  \item \textbf{Testing page size fairly.} Changing the page size alone changes the
        amount of memory. Generating addresses as byte positions and holding memory
        constant isolated the real effect.
  \item \textbf{Using one random seed.} The first LRU and FIFO comparison on uniform
        references suggested that FIFO was better. Repeating over 10 seeds showed that
        the policies were equal, and that the gap was noise.
\end{itemize}

\subsection{Enhancements Beyond the Requirements}
\begin{itemize}
  \item \code{--compare} prints side-by-side tables from a single seeded workload.
  \item \code{--seed} makes every run repeatable.
  \item A TLB size of 0 turns the TLB off, which gives a measured baseline.
  \item Second-level tables are allocated on first use, and the simulator reports
        table memory for two-level and single-level designs.
  \item Stress workloads have adjustable locality with \code{--locality}.
  \item Exceptions form a hierarchy, and a page fault is handled as a trap and retry.
  \item A regression suite runs under AddressSanitizer and UndefinedBehaviorSanitizer.
\end{itemize}

\subsection{Summary of Results}
LRU beat FIFO only when the workload had locality. A TLB cut average latency by up to
47\% at seed 42 without affecting page faults, with most of the gain coming from the
first few entries. Larger pages improved TLB reach and reduced faults slightly, at a
cost in internal fragmentation that the simulator does not measure. Two-level paging
saves table memory only in sparse address spaces, and its extra latency is acceptable
only because of the TLB. Latency values come from a simple cost model, so the
comparisons between configurations are more meaningful than the absolute numbers.

%=====================================================================
\section{References}
\begin{enumerate}
  \item R. H. Arpaci-Dusseau and A. C. Arpaci-Dusseau, \textit{Operating Systems: Three
        Easy Pieces}, chapters on TLBs, advanced page tables and swapping policies.
        \url{https://pages.cs.wisc.edu/~remzi/OSTEP/}
  \item P. J. Denning, ``The Working Set Model for Program Behavior,''
        \textit{Communications of the ACM}, vol. 11, no. 5, 1968.
  \item W. W. Chu and H. Opderbeck, ``The Page Fault Frequency Replacement Algorithm,''
        \textit{Proceedings of the AFIPS Fall Joint Computer Conference}, 1972.
  \item CS375 Lab 9 handout and starter code.
\end{enumerate}

\end{document}
